\section{Actionable Protocol: Five Injectable Templates}
\label{sec:protocol}

We distill the framework into five templates (T1--T5), each with a trigger
condition and an action. Each template is a short instruction that can be
injected into an agent prompt as a thinking habit. We position them as an
\emph{agent-prompt injection protocol}: the claim is that they are
\emph{injectable and operable} (well-defined triggers, executable actions,
composable). We do \emph{not} claim effectiveness; that is pending independent
experiments.

\begin{enumerate}[leftmargin=2em,itemsep=2pt]
\item \textbf{T1 Cost model.} \emph{Trigger:} a concrete solution has been
found, or ``is $X$ optimal?'' is asked. \emph{Action:} write down the cost model
(what is being costed, under what assumptions) and the unified form; verify the
original solution is a special case; if it cannot be derived, name the missing
assumption.
\item \textbf{T2 Unified coverage.} \emph{Trigger:} facing several seemingly
unrelated methods/systems. \emph{Action:} find one formula or analogy covering
all of them (compression rate / information gain per cost / channel bandwidth),
and write out each term's degeneracy conditions.
\item \textbf{T3 Quantitative landing.} \emph{Trigger:} a qualitative conclusion
is given. \emph{Action:} immediately quantify it into a testable number
(success-rate expression / bandwidth lower bound / expected number of tests); if
it cannot be quantified, mark it ``needs measurement'' and design the cheapest
probe.
\item \textbf{T4 Analogy audit.} \emph{Trigger:} about to borrow a historical
case / analogy / prior. \emph{Action:} check dimension-by-dimension the
condition difference $\Delta$ and the conclusion's sensitivity $s_k$; if
sensitive and mismatched, do not borrow; attach an $\varepsilon\cdot$uniform
floor to prevent confident errors.
\item \textbf{T5 Criterion bottleneck.} \emph{Trigger:} a framework has taken
shape. \emph{Action:} find the ``criterion bottleneck''---which step in the loop
lacks cheap, reliable ground truth; that step is the real bottleneck, not the
architecture. Prefer the remedy ``make the world closed'' (synthetic domains /
executable verification).
\end{enumerate}

\paragraph{Positioning.}
The five templates are the operational layer of the framework: T1--T2 are the
``climb one level'' move (concrete $\to$ cost model $\to$ unified form, original
solution as special case); T3 is predictive power (conclusions land on testable
numbers); T4 is the robustness discipline of leading-term anchoring applied to
borrowed priors; T5 is the fitness-landscape move (locate where ground truth is
missing). Their value proposition is \emph{operability}: each has a crisp
trigger and a mechanical action, so they can be injected into a prompt and
audited for whether they fired. Whether they improve downstream task performance
is an empirical question we leave open.
